\newpage
\phantomsection
\label{prf:metric-point-functions-identify-points}
\LRAProofFor{thm:metric-point-functions-identify-points}

\begin{remark*}[Return]
\hyperref[thm:metric-point-functions-identify-points]{Return to Theorem}
\end{remark*}

\begin{theorem*}[Point Functions Identify Points]
Let \((X,d)\) be a metric space.  The map
\[
  X\to\delta(X),
  \qquad
  z\mapsto\delta_z
\]
is a bijection.
\end{theorem*}

\begin{proof}[Professional Standard Proof]
\LRAProofBodyStart
Let
\[
  F:X\to\delta(X),
  \qquad
  F(z)=\delta_z.
\]
The map \(F\) is surjective because \(\delta(X)\) is defined to be the set of
all functions \(\delta_z\) with \(z\in X\).  Thus every element of
\(\delta(X)\) has the form \(F(z)\) for some \(z\in X\).

It remains to prove injectivity.  Suppose \(z_1,z_2\in X\) and
\[
  F(z_1)=F(z_2).
\]
Then \(\delta_{z_1}=\delta_{z_2}\) as functions.  Evaluating this equality at
\(z_1\) gives
\[
  \delta_{z_1}(z_1)=\delta_{z_2}(z_1).
\]
By the definition of point function,
\[
  d(z_1,z_1)=d(z_2,z_1).
\]
The left side is \(0\), so \(d(z_2,z_1)=0\).  By the identity of
indiscernibles for a metric, \(z_2=z_1\), and hence \(z_1=z_2\).  Therefore
\(F\) is injective.

Since \(F\) is both surjective and injective, it is bijective.
\end{proof}

\begin{proof}[Detailed Learning Proof]
\LRAProofBodyStart
We need to prove that the assignment
\[
  F:X\to\delta(X),
  \qquad
  F(z)=\delta_z,
\]
is a bijection.  This means proving two things: every element of the codomain
is hit by \(F\), and no two different points of \(X\) are sent to the same
point function.

First consider surjectivity.  The codomain is
\[
  \delta(X)=\{\,\delta_z:z\in X\,\}.
\]
So if \(u\in\delta(X)\), then by definition there is some \(z\in X\) such
that \(u=\delta_z\).  But \(F(z)=\delta_z\), so \(F(z)=u\).  Hence \(F\) is
surjective onto \(\delta(X)\).

Now consider injectivity.  Suppose two points \(z_1,z_2\in X\) satisfy
\[
  F(z_1)=F(z_2).
\]
Since \(F(z)=\delta_z\), this says
\[
  \delta_{z_1}=\delta_{z_2}.
\]
These are functions on \(X\), so they have the same value at every input.
Choose the input \(x=z_1\).  Then
\[
  \delta_{z_1}(z_1)=\delta_{z_2}(z_1).
\]
Using the definition \(\delta_z(x)=d(z,x)\), this becomes
\[
  d(z_1,z_1)=d(z_2,z_1).
\]
A metric has zero self-distance, so \(d(z_1,z_1)=0\).  Therefore
\[
  d(z_2,z_1)=0.
\]
The identity of indiscernibles says that distance \(0\) occurs only between
equal points.  Thus \(z_2=z_1\), and so \(z_1=z_2\).  This proves that \(F\)
is injective.

The map \(F\) is surjective because \(\delta(X)\) is exactly its image, and it
is injective because equality of distance profiles forces equality of the
points that produced them.  Therefore \(F\) is bijective.
\end{proof}

\begin{remark*}[Proof structure]
Surjectivity is immediate from the definition of \(\delta(X)\).  Injectivity
uses equality of functions at a strategically chosen input: if
\(\delta_{z_1}=\delta_{z_2}\), evaluate both sides at \(z_1\).  The metric
axiom \(d(z_1,z_1)=0\), together with identity of indiscernibles, forces
\(z_2=z_1\).
\end{remark*}

\begin{dependencies}
\begin{itemize}
  \item \hyperref[def:metric-space]{Definition: Metric Space}
  \item \hyperref[def:metric-point-function]{Definition: Point Function}
\end{itemize}
\end{dependencies}

\clearpage
